\documentclass[11pt]{article}
\usepackage[T1]{fontenc}
\usepackage{lmodern,microtype,amsmath,amssymb,amsthm,mathtools}
\usepackage[letterpaper,margin=1in]{geometry}
\usepackage[dvipsnames]{xcolor}
\usepackage{fancyhdr}
\usepackage[colorlinks=true,linkcolor=MidnightBlue,citecolor=MidnightBlue,urlcolor=MidnightBlue]{hyperref}
\pagestyle{fancy}\fancyhf{}\lhead{Factorial samples of prime-power sums}\rhead{Research note}\cfoot{\thepage}
\setlength{\headheight}{14pt}
\newtheorem{theorem}{Theorem}
\newtheorem{lemma}[theorem]{Lemma}
\newtheorem{corollary}[theorem]{Corollary}
\newcommand{\R}{\mathbb R}
\newcommand{\C}{\mathbb C}
\newcommand{\e}{\mathrm e}
\newcommand{\dd}{\mathrm d}
\DeclareMathOperator{\Rea}{Re}
\title{Factorial samples of prime-power sums\\[2pt]\large An oscillating autocorrelation determines a primitive Dirichlet character}
\author{}
\date{26 November 2026}
\begin{document}
\maketitle\thispagestyle{plain}
\begin{abstract}
At the factorial endpoints we sample a smoothed, centered sum of prime
powers twisted by a primitive Dirichlet character. Assuming its
critical-line zero theorem, we prove that its mean autocorrelation at
every positive integer lag oscillates indefinitely. We identify the
entire oscillating term. Its logarithmic Fourier coefficients recover
every zero multiplicity, and hence determine the character. We also
obtain an exact distance formula for pairs of characters and transfer
the joint laws at fixed multiplicative shifts to factorial samples.
The proof allows multiple zeros and arbitrary rational relations among
their ordinates.
\end{abstract}

\section{The arithmetic measurements}
Write $\Lambda(p^k)=\log p$ for primes $p$ and $k\ge1$, and
$\Lambda(n)=0$ otherwise. All logarithms are natural. For a primitive
Dirichlet character $\chi$ of conductor $q$, set $d_\chi=1$ when $q=1$
(the zeta function), and $d_\chi=0$ otherwise. For $x>0$ put
\begin{equation}\label{eq:B}
 B_\chi(x)=\frac{1}{\sqrt x}
 \left(\sum_{n\le x}\Lambda(n)\chi(n)(1-n/x)-d_\chi x/2\right).
\end{equation}
The arithmetic weights vanish at the endpoint. Let
$m_\chi(t)$ be the multiplicity of $1/2+it$ as a zero, with value
zero elsewhere, and set
\[
 c_\chi(t)=-\frac{m_\chi(t)}{(1/2+it)(3/2+it)}.
\]
Sums over $t$ run over distinct real ordinates with nonzero summands.
For $h\in\mathbb Z_{\ge1}$, define the running lag-$h$ autocorrelation
\[
 C_{\chi,h}(N)=\frac1N\sum_{n=1}^{N}B_\chi((n+h)!)\overline{B_\chi(n!)}.
\]

\begin{theorem}\label{th:lag}
Suppose every nontrivial zero of $L(s,\chi)$ has real part $1/2$. Then
\begin{equation}\label{eq:lag}
 C_{\chi,h}(N)=
 \sum_t\frac{|c_\chi(t)|^2}{1+iht}\,\e^{iht\log N}+o(1).
\end{equation}
The series converges absolutely and uniformly as a function of $\log N$.
The sequence $C_{\chi,h}(N)$ has no limit. For every real $v$,
\begin{equation}\label{eq:outer}
 \lim_{X\to\infty}\frac{1}{\log X}\sum_{N\le X}
 \frac{C_{\chi,h}(N)}{N^{1+ihv}}
 =\frac{|c_\chi(v)|^2}{1+ihv}.
\end{equation}
In particular, the running averages at $h=1$ determine $\chi$.
\end{theorem}

\newpage
\section{Two analytic ingredients}
\begin{lemma}[Smoothed explicit formula]\label{lem:explicit}
Under the hypothesis of Theorem~\ref{th:lag}, for $x\ge2$,
\begin{equation}\label{eq:explicit}
 B_\chi(x)=\sum_t c_\chi(t)\e^{it\log x}
 +O_\chi\big((1+\log x)/\sqrt x\big),
 \qquad \sum_t |c_\chi(t)|<\infty .
\end{equation}
\end{lemma}

\begin{proof}
Mellin inversion of $(1-u)_+$, whose transform for $\Rea s>0$
is $1/[s(s+1)]$, gives the exact unnormalized prime sum as
\[
 \frac1{2\pi i}\int_{2-i\infty}^{2+i\infty}
 -\frac{L'}{L}(s,\chi)\frac{x^s}{s(s+1)}\,\dd s .
\]
The series for $-L'/L$ converges absolutely on this line and the kernel
is integrable. Shift to $\Rea s=-1/4$. We give the needed estimates to
specify the treatment of the infinite contour. On the new line the
functional equation, Stirling's formula and the absolutely convergent
logarithmic derivative at $1-s$ give
$L'/L\ll_\chi\log(2+|\Im s|)$. Thus its integral is $O_\chi(x^{-1/4})$.
On the horizontal edges one may choose positive and negative heights
tending to infinity with $L'/L\ll_\chi\log^2(2+|\Im s|)$ throughout
the strip. Indeed the classical local logarithmic derivative formula
on this fixed strip subtracts the simple fractions $1/(s-\rho)$ for
$|\Im\rho-\Im s|\le1$ with an $O_\chi(\log(2+|\Im s|))$ error.
Each unit interval contains $O_\chi(\log(2+|\Im s|))$ ordinates
counted with multiplicity. Removing intervals of radius
$\varepsilon_\chi/\log(2+|\Im s|)$ around the nearby ordinates
therefore leaves an admissible height for sufficiently small
$\varepsilon_\chi>0$. These standard fixed-conductor estimates follow
from the completed Hadamard product and Stirling's formula
\cite{davenport,bennett}. The quadratic kernel makes both horizontal
integrals tend to zero for fixed $x$.

The pole of $L$ at $1$, when present, contributes $+x/2$.
A nontrivial zero $\rho$ of multiplicity $m$ contributes
$-m x^\rho/[\rho(\rho+1)]$. At $s=0$, if
$-L'/L=-r/s+b+O(s)$ with $r=\operatorname{ord}_0 L$, the residue is
$b-r(\log x-1)$. There are no other poles between the two lines.
Division by $\sqrt x$ gives~\eqref{eq:explicit}. The zero count up to
height $T$ is $O_\chi(T\log(2+T))$; partial summation against
$(1+|t|)^{-2}$ gives the required absolute convergence.
\end{proof}

\begin{lemma}[Factorial phases]\label{lem:phase}
For any real $u\ne0$ and $N\ge2$,
\[
 \left|\sum_{n=1}^{N}\e^{iu\log(n!)}\right|
 \ll \sqrt N\,\big(\sqrt{|u|}+1/\sqrt{|u|}\big).
\]
For each fixed $v\in\R$, insertion of the additional weight $n^{iv}$
leaves this bound valid with a constant depending on $v$.
\end{lemma}
\begin{proof}
On $M\le x\le2M$ the second derivative of
$f(x)=u\log\Gamma(x+1)/(2\pi)$ has constant sign and absolute size
$\asymp |u|/M$, by
$(\log\Gamma)''(x+1)=\sum_{j\ge0}(x+1+j)^{-2}$.
Van der Corput's second-derivative estimate bounds the exponential sum
on any subinterval by $O(M\sqrt{|u|/M}+\sqrt{M/|u|})$.
Summing dyadic blocks proves the first statement. Partial summation
and $|(x^{iv})'|=|v|/x$ prove the second. Factorial equidistribution
also appears in Diaconis's work on leading digits \cite{diaconis}.
\end{proof}

\newpage
\section{What the lag averages remember}
\begin{proof}[Proof of Theorem~\ref{th:lag}, except character determination]
By Lemma~\ref{lem:explicit} we can replace each $B$ by its uniformly
convergent trigonometric series. The replacement error vanishes in the
mean because both factors are bounded and their errors tend to zero.
First truncate both series to finite frequency sets. For fixed
$t,u\in\R$, the mean of the corresponding pair of modes involves
\begin{align*}
 \e^{it\log((n+h)!)-iu\log(n!)}
 &=\e^{i(t-u)\log(n!)}\,
   \exp\!\left(it\sum_{j=1}^h\log(n+j)\right)\\
 &=\e^{i(t-u)\log(n!)}\,n^{iht}(1+O_{h,t}(1/n)).
\end{align*}
When $t\ne u$ the mean tends to zero by the weighted version of
Lemma~\ref{lem:phase}. For $t=u$, integral comparison gives
\[
 \frac1N\sum_{n=1}^N n^{iht}=\frac{N^{iht}}{1+iht}
 + O_{h,t}\big((1+\log N)/N\big).
\]
In fact the sum differs from $\int_1^N x^{iht}\,\dd x$ by
$O(1+|ht|\log N)$. This proves~\eqref{eq:lag} for finite
truncations. The double-series tail is uniformly small by
$\sum_t |c_\chi(t)|<\infty$, and the limiting diagonal series
also converges absolutely. Taking the truncations to infinity proves
the asserted $o(1)$.

For real $w$ the logarithmic mean of $N^{iw}$ equals $1$ for $w=0$
and $0$ otherwise. Indeed the difference between
$\sum_{N\le X}N^{-1+iw}$ and $\int_1^X x^{-1+iw}\,\dd x$ is
$O_w(1)$, by integrability of the derivative. Insert~\eqref{eq:lag}
in the left side of~\eqref{eq:outer}. The error has logarithmic mean
zero, and absolute convergence permits termwise passage to the limit.
This gives~\eqref{eq:outer}, including at $v=0$.

Every primitive Dirichlet $L$-function has nontrivial zeros of
arbitrarily large height by its standard zero-counting asymptotic
\cite{bennett,davenport}. Choose a nonzero ordinate $v$. If
$C_{\chi,h}(N)$ tended to a constant, its logarithmic coefficient at
$hv$ would be zero by the preceding elementary mean formula. Equation
\eqref{eq:outer} gives a nonzero coefficient there. Finally,
multiplication of~\eqref{eq:outer} by $1+ihv$ yields $|c_\chi(v)|^2$.
The formula
\[
 m_\chi(v)=\sqrt{(1/4+v^2)(9/4+v^2)}\,|c_\chi(v)|
\]
therefore recovers every multiplicity, including zero for an absent
ordinate. Character determination follows from the next lemma.
\end{proof}

\begin{lemma}[Determination by the zero divisor]\label{lem:id}
Two primitive Dirichlet characters with the same multiset of nontrivial
zeros are equal.
\end{lemma}
\begin{proof}
Set $a_\chi\in\{0,1\}$ by $\chi(-1)=(-1)^{a_\chi}$ and put
\[
 F_\chi(s)=[s(s-1)]^{d_\chi}
 (q_\chi/\pi)^{(s+a_\chi)/2}
 \Gamma\big((s+a_\chi)/2\big)L(s,\chi).
\]
This is entire of order one with exactly the stated zeros, including
for $q_\chi=1$. Hadamard factorization gives
$F_\chi/F_\psi=\e^{A+Bs}$ when the multisets coincide.
Along $s=\sigma\to+\infty$, the Dirichlet series both tend to $1$.
Stirling's ratio formula yields
\[
 \log|F_\chi/F_\psi|(\sigma)
 =\tfrac12\sigma\log(q_\chi/q_\psi)
 +\big(2(d_\chi-d_\psi)+(a_\chi-a_\psi)/2\big)\log\sigma+O(1).
\]
The exact affine dependence on $\sigma$ forces the logarithmic
coefficient to vanish. Since $a_\chi,a_\psi,d_\chi,d_\psi\in\{0,1\}$,
both differences vanish. Removing the common polynomial and gamma
factors leaves $L(s,\chi)/L(s,\psi)$ exponential in an affine
function of $s$. Along the positive real axis it tends to one,
so it is identically one. Uniqueness of coefficients of an absolutely
convergent Dirichlet series gives $\chi(n)=\psi(n)$ for every $n$.
Their common arithmetic function has a unique primitive conductor.
\end{proof}

\section{Direct recovery and joint laws}
\begin{corollary}\label{cor:distance}
Assume the critical-line zero theorem for the primitive characters
under consideration. For every $v\in\R$,
\begin{equation}\label{eq:direct}
 \lim_{N\to\infty}\frac1N\sum_{n\le N} B_\chi(n!)\,(n!)^{-iv}=c_\chi(v).
\end{equation}
For two primitive characters $\chi,\psi$,
\begin{equation}\label{eq:metric}
 \lim_{N\to\infty}\frac1N\sum_{n\le N}|B_\chi(n!)-B_\psi(n!)|^2
 =\sum_t\frac{|m_\chi(t)-m_\psi(t)|^2}{(1/4+t^2)(9/4+t^2)}.
\end{equation}
The square root is a metric on primitive characters.
\end{corollary}
\begin{proof}
Insert~\eqref{eq:explicit} in~\eqref{eq:direct}. Every frequency except
$v$ vanishes in the mean by Lemma~\ref{lem:phase}. For~\eqref{eq:metric}
expand the squared absolute value using frequencies from both
characters. Absolute convergence justifies termwise means, and the
distinct-frequency terms vanish. The right side is a squared weighted
$\ell^2$ distance. Zero distance gives equal zero multisets under
the critical-line hypotheses, so Lemma~\ref{lem:id} gives equality.
\end{proof}

\begin{corollary}[Fixed multiplicative shifts]\label{cor:joint}
Fix finitely many primitive characters $\chi_j$ satisfying the
critical-line statement and $b_j\in\R$, $1\le j\le r$.
Set $V(y)=(B_{\chi_j}(\e^{y+b_j}))_{j=1}^r$. For every continuous
$H:\C^r\to\C$, both limits exist and satisfy
\[
 \lim_{N\to\infty}\frac1N\sum_{n\le N}H(V(\log(n!)))
 =\lim_{Y\to\infty}\frac1Y\int_0^Y H(V(y))\,\dd y.
\]
\end{corollary}
\begin{proof}
The inputs stay in a compact box for $y\ge0$. Replacing each
coordinate by the bounded trigonometric series
$\sum_t c_{\chi_j}(t)\e^{it(y+b_j)}$ introduces a difference tending
to zero, also after composition with $H$ by uniform continuity on
a common compact box. Approximate the new $H$ of the vector
uniformly by finite trigonometric polynomials: first approximate $H$
on the box by a polynomial in the coordinates and their conjugates
(Stone--Weierstrass), then truncate each coordinate series. Each
trigonometric polynomial has the same mean by both procedures,
because a nonzero-frequency mode has continuous mean zero and
Lemma~\ref{lem:phase} gives factorial mean zero. Uniform
approximation transfers existence and equality of limits.
\end{proof}

\paragraph{The role of primitivity.}
Suppose a character modulo $Q$ is induced by a primitive character
$\chi^*$. Use the same definition~\eqref{eq:B}, taking the main term
from $\chi^*$. Only powers of primes dividing $Q$ can give discrepancies
between the two weighted sums. Each such prime contributes at most
$\log x$ in absolute value for $x\ge2$. Consequently
\[
 |B_\chi(x)-B_{\chi^*}(x)|\le
 \omega(Q)\log(x)/\sqrt x,
\]
where $\omega(Q)$ counts distinct prime divisors. Thus the experiments
recover precisely the inducing primitive character if imprimitive
inputs are admitted. For example all principal characters have the
same limiting measurements as the zeta function.

\newpage
\section{Interpretation and provenance}
At two fixed multiplicative shifts the logarithmic time separation is
constant. Between two factorial endpoints $h$ sample indices apart,
this separation instead equals $h\log n+O_h(1/n)$. On the diagonal
of the double zero sum this produces $n^{iht}$. Its arithmetic mean
retains the oscillation at the upper endpoint. This accounts for the
different conclusions of Theorem~\ref{th:lag} and
Corollary~\ref{cor:joint}.

The input on zero locations is the statement about primitive
Dirichlet characters in the linked Lean development~\cite{lean}.
Its conductor-one case is RH. The linked overview~\cite{wiki} also
describes the explicit-formula connection between zeta zeros and
prime counts. Here the triangular weight in~\eqref{eq:B} adds a
second denominator to the Mellin kernel. That extra decay is what
permits uniform control on a set of scales with very large gaps.

Factorial phase cancellation, explicit formulae and determination by
a completed zero divisor are classical. For related continuous
limiting-distribution theorems, see Akbary, Ng and Shahabi
\cite{akbary}. The contribution proposed here is the factorial
lag asymptotic~\eqref{eq:lag}, its second-stage inversion
\eqref{eq:outer}, and their arithmetic application. No independent
priority determination is claimed. These are exact infinite-data
identities. The bounds in Lemma~\ref{lem:phase} deteriorate when
two frequencies approach one another. An effective sample size for
separating characters of bounded conductor would require quantitative
control of spectral tails and of the low-lying distinguishing ordinates.

\begin{thebibliography}{9}\small
\bibitem{lean} L.~Sernau et al., \emph{GRH for Dirichlet $L$-functions, in
Lean 4}, source repository, accessed 26 November 2026.
\url{https://github.com/lsernau/grh-lean}.
\bibitem{wiki} Wikipedia contributors, \emph{Riemann hypothesis}, especially
the discussion of prime distribution. Accessed 26 November 2026.
\url{https://en.wikipedia.org/wiki/Riemann_hypothesis}.
\bibitem{davenport} H.~Davenport, \emph{Multiplicative Number Theory},
3rd ed., revised by H.~L.~Montgomery, Springer, 2000.
\href{https://doi.org/10.1007/978-1-4757-5927-3}{doi:10.1007/978-1-4757-5927-3}.
\bibitem{bennett} M.~A.~Bennett, G.~Martin, K.~O'Bryant and A.~Rechnitzer,
Counting zeros of Dirichlet $L$-functions,
\emph{Math. Comp.} \textbf{90} (2021), 1455--1482.
\href{https://doi.org/10.1090/mcom/3599}{doi:10.1090/mcom/3599}.
\bibitem{diaconis} P.~Diaconis, The distribution of leading digits and
uniform distribution mod 1, \emph{Ann. Probab.} \textbf{5} (1977), 72--81.
\href{https://doi.org/10.1214/aop/1176995891}{doi:10.1214/aop/1176995891}.
\bibitem{akbary} A.~Akbary, N.~Ng and M.~Shahabi,
Limiting distributions of the classical error terms of prime number theory,
\emph{Q. J. Math.} \textbf{65} (2014), 743--780.
\href{https://www.cs.uleth.ca/~akbary/limitingdistributionFINALQJM.pdf}{Authors' paper}.
\end{thebibliography}
\end{document}
